%% =====================================================================================
%% P-3 — Setting, formalisation, panel, protocol.  Muc LaTeX DOC LAP.
%% Nhan `sec:methods` LA CHO DUA CUA MUC §6 (`\ref{sec:methods}` trong 06_bestresponse.tex).
%% Doi nhan nay thi phai sua ca o do.
%%
%% MENH DE 1 DA DUOC KIEM BANG VET CAN, khong phai chi chung minh tren giay:
%%     python -m pytest crsd/tests/test_paper_equilibrium.py
%% Test do duyet moi profile doi xung + 6.000 profile bat doi xung tai 11 muc p va doi
%% chieu voi dac trung phat bieu o day. SUA MENH DE THI PHAI SUA TEST.
%%
%% ⚠️ p = 1 LA NGOAI LE THAT, DA KIEM: khi p = 1, moi profile truot target deu la can bang
%% (yeu) vi payoff bang 0 du lam gi. Vi the menh de phat bieu cho p < 1. Bo chu "p<1" di
%% la sai, va la thu reviewer kiem duoc bang mot dong code.
%%
%% NGAN SACH TRANG: 1,7 cot / 16 cot. Khong float.
%% =====================================================================================

\section{The game, its equilibria, and the panel}
\label{sec:methods}

\subsection{The game}

The \PanelNplayersWord{} players each hold an endowment of $e=\PanelEndowment$ units. Over
$R=\PanelNrounds$ rounds every player privately chooses a contribution
$a_i^t\in\{0,2,4\}$, and after each round every player sees what all six chose in every
round so far. Writing $x_i=\sum_t a_i^t$ for a player's total and $X=\sum_i x_i$ for the
group's, the group succeeds if $X\ge T=\PanelTarget$. A player's cash is then what it kept:
\begin{equation}
u_i(x) \;=\;
\begin{cases}
e-x_i, & X\ge T,\\[2pt]
0 \text{ w.p.\ } p,\quad e-x_i \text{ w.p.\ } 1-p, & X<T .
\end{cases}
\label{eq:payoff}
\end{equation}
Contributions are never refunded, the round count is stated in the prompt, and
$R\cdot\max A=e$, so no player can contribute more than it holds. The risk $p$ is the only
quantity we vary in Section~\ref{sec:riskgrid}. Since \eqref{eq:payoff} depends on the
action profile only through the totals $x$, the Nash equilibrium \emph{outcomes} of the
$R$-round game are those of the reduced game in which each player picks
$x_i\in\{0,2,\dots,e\}$, and we characterise them there.

\begin{proposition}
\label{prop:eq}
Fix $p<1$. A profile $x$ is a Nash equilibrium if and only if either $x=0$, or $\sum_i
x_i=T$ and $x_i\le p\,e$ for every $i$. Hence an equilibrium that reaches the target exists
if and only if $p\ge p^{*}$ with $p^{*}=T/(Ne)=\PanelPstar$, and at $p=p^{*}$ the only such
equilibrium is the equal split $x_i=T/N=\PanelFairTotal$. The welfare-maximising outcome is
$x=0$ for $p<p^{*}$ and $\sum_i x_i=T$ for $p>p^{*}$, so the per-capita payoff of the best
outcome is $\max\{(1-p)e,\;T/N\}$.
\end{proposition}

\noindent\emph{Proof sketch.} If $X<T$ and $x_i>0$, then moving to $x_i=0$ leaves the
target missed and raises the expectation from $(1-p)(e-x_i)$ to $(1-p)e$, strictly for
$p<1$; so an equilibrium that misses the target has $x=0$, and $x=0$ is one because a lone
player can contribute at most $e<T$. If $X>T$ then $X\ge T+2$, and any $i$ with $x_i\ge2$
sheds $2$, still reaches the target and gains $2$; so an equilibrium that reaches it has
$X=T$ exactly. Given $X=T$, every profitable-looking deviation reduces $x_i$ and therefore
misses the target, and the best of them is $x_i=0$, worth $(1-p)e$ against $e-x_i$ for
staying; the two compare as $x_i\le p\,e$. Summing that bound over $N$ players against
$\sum_i x_i=T$ gives $p\ge T/(Ne)$, with equality forcing $x_i=T/N$ for all $i$. For
welfare, $\sum_i u_i$ equals $Ne-X$ when $X\ge T$ and $(1-p)(Ne-X)$ otherwise, so it is
maximised at $X=T$ or at $X=0$, and $Ne-T\ge(1-p)Ne$ reduces to $p\ge T/(Ne)$
again.\hfill\crgqed

Three consequences carry the rest of the paper. Welfare-optimal play and best-equilibrium
play switch at the \emph{same} threshold, so the benchmark
$\mathrm{opt}(p)=\max\{(1-p)e,T/N\}$ used throughout needs no choice between the two.
Below $p^{*}$ no equilibrium reaches the target, which is to say that a group that reaches
it there is not merely being generous but is playing an outcome no player should sustain.
And at $p=0$ the two branches of \eqref{eq:payoff} coincide, so $u_i=e-x_i$ whatever the
group does and contributing is strictly dominated under \emph{any} increasing utility,
not only under the risk-neutral one that $\mathrm{opt}$ assumes. The $p=0$ column is
therefore the one place in the paper where no assumption about risk attitude is doing any
work, and we lean on it accordingly.

\subsection{Panel and protocol}

%% CHO DUY NHAT trong bai in slug day du. Moi cho khac dung ten ngan trong ngoac o day —
%% xem `FULL` trong panel_analysis.py, no lay thang tu `MODELS` cua e3a_analysis.py nen ca
%% hai muc ket qua goi model bang cung mot ten.
We run \PanelNmodelsWord{} production models, one per provider and each at the inexpensive
tier of its family, so that no result is an artefact of a single vendor's post-training:
\textsc{claude-haiku-4.5} (Haiku), \textsc{gemini-3.5-flash-lite} (Flash-Lite),
\textsc{gpt-5.6-luna} (Luna), \textsc{qwen3-235b-a22b-instruct} (Qwen) and
\textsc{grok-4.20-non-rea\-son\-ing} (Grok), named by their short forms from here on. Every
experiment runs every model on exactly the same number of games; price never determines
sample size. The main panel is in English, at temperature $0.7$, with every seat holding a
neutral persona. E6 adds two prompt paraphrases and a temperature-zero arm at the endpoint
risk cells $p\in\{0.1,0.9\}$, with \EsixReps{} games per model, arm, and risk level.

Each round a seat receives the rules, the complete per-seat history in fixed positions, and
a request for one contribution. The running group total is \emph{not} printed, so a seat
that wants it must add the history up; we keep it that way because printing it changes
behaviour, and because Section~\ref{sec:bestresponse} needs a measure that uses only what
the prompt displays. A reply is parsed from a \texttt{CONTRIBUTION:}\ line. Replies that
reach the output cap before that line are retried at a quadrupled cap rather than parsed,
because a truncated chain of arithmetic otherwise yields a plausible number that was never
a decision. Across the main panel and diagnostic experiments there are $\PanelNgamesAll$ games and
$\PanelNseatRoundActions$ seat--round actions. The self-play/control panel contributes
$\PanelNdecisionsPanel$ model-seat decisions; the scripted-opponent experiment contributes
one LLM decision per round, and the mixed-population experiment contributes six. In total,
the study contains $\PanelNllmDecisionsAll$ LLM decisions; all reported games have zero
parse failures and zero truncations.

Cells hold $\PanelNpercell$ games each, indexed by a repetition number that also seeds the
catastrophe lottery, so the same repetition faces the same draw at every risk level.
Those common random numbers are what make the paired comparisons across $p$ in
Section~\ref{sec:riskgrid} worth making. The unit of analysis is the game, never the seat:
six seats in one game read each other's history and are not independent. Every interval is
a percentile bootstrap over games with $\PanelNboot$ resamples and every $P$-value a
permutation test with $\PanelNpermTex$ resamples, both under a fixed seed.
